% Part: normal-modal-logic
% Chapter: axioms-systems
% Section: introduction

\documentclass[../../../include/open-logic-section]{subfiles}

\begin{document}

\olfileid[tr]{nml}{axs}{int}

\olsection{Giriş}

Temel kiplik dili için kiplik modelleri cinsinden bir semantiğimiz vardır.
Ayrıca !!a{formula} için geçerli olma---bütün modellerin bütün dünyalarında
doğru olma---ya da belirli bir modeller veya çerçeveler sınıfına göre
geçerli olma---sınıftaki veya çerçeveye dayanan bütün modellerin bütün
dünyalarında doğru olma---kavramına sahibiz. Mantık genellikle bu semantik
geçerlilik nitelemelerini ispat kuramsal bir !!{derivability} kavramıyla
ilişkilendirir. Amaç, bir sistem içinde öyle bir !!{derivability} kavramı
tanımlamaktır ki !!a{formula} ancak ve ancak geçerliyse !!{derivable} olsun.

En basit ve tarihsel olarak en eski !!{derivation} sistemleri, Hilbert tipi
ya da aksiyomatik !!{derivation} sistemleri denen sistemlerdir. Birçok modal
mantık için Hilbert tipi !!{derivation} sistemleri kurmak görece kolaydır:
metateorik incelemenin nesneleri olarak basittirler (örneğin sağlamlık ve
tamlıklarını ispatlamak kolaydır). Buna karşılık bu sistemlerde !!{formula}s
ispatlamak, örneğin doğal tümdengelim sistemlerinde ispatlamaktan çok daha
zordur.

Hilbert tipi !!{derivation} sistemlerinde !!a{derivation}, bir
!!a{formula} için belirli aksiyomlardan başlayıp birkaç çıkarım kuralı
aracılığıyla söz konusu !!{formula} ile biten bir !!{formula}s dizisidir.
!!{derivation} sisteminin
semantik ile uyuşmasını istediğimizden, !!{derivable} formüller kümesinin
bütün modellerde (ya da bütün aksiyomların doğru olduğu bütün modellerde)
doğru olduğunu güvence altına almalıyız. Önce bunun işlemesi için gerekli
olan bazı modal mantık özelliklerini, yani ``normal'' modal mantıkları
ayıracağız. Normal modal mantıklar için yalnızca iki çıkarım kuralını
varsaymak gerekir: modus ponens ve zorunlulaştırma. Aksiyom olarak bütün
totolojileri (onların yerine koyma örneklerini) ve ele aldığımız modal mantığa
bağlı olarak bazı modal aksiyomları alırız. Yalnızca bütün modeller sınıfıyla
ilgilensek bile $\Ax{K}$'nin\iftag{notprvDiamond}{}{ ve~$\Ax{Dual}$'ın}
bütün yerine koyma örneklerini de aksiyom saymalıyız. Yalnızca bunlar,
en küçük normal modal mantık $\Log K$'yi üretir.

\begin{defn}
\emph{Modus ponens} kuralı şu çıkarım şemasıdır:
\begin{prooftree}
\AxiomC{$!A$}
\AxiomC{$!A \lif !B$}
\RightLabel{\MP}
\BinaryInfC{$!B$}
\end{prooftree}
$!C\ident !A\lif !B$ ise, !!a{formula}~$!B$'nin $!A$ ve $!C$
!!{formula}s kullanılarak modus ponens ile \emph{çıktığını} söyleriz.
\end{defn}

\begin{defn}
\emph{Zorunlulaştırma} kuralı şu çıkarım şemasıdır:
\begin{prooftree}
\AxiomC{$!A$}
\RightLabel{\Nec}
\UnaryInfC{$\Box !A$}
\end{prooftree}
$!B\ident\Box !A$ ise !!{formula}~$!B$'nin !!{formula}~$!A$'dan
zorunlulaştırma ile çıktığını söyleriz.
\end{defn}

\begin{defn}
Bir $\Sigma$ aksiyomlar kümesinden bir \emph{!!{derivation}}, her $!B_i$'nin
aşağıdakilerden biri olduğu bir $!B_1,!B_2,\dots,!B_n$ !!{formula}s dizisidir:
\begin{enumerate}
\item bir totolojinin yerine koyma örneği;
\item $\Sigma$ içinde bulunan !!a{formula} için bir yerine koyma örneği;
\item $j,k<i$ olmak üzere önceki $!B_j$, $!B_k$ !!{formula}s kullanılarak modus
  ponens ile çıkan bir formül;
\item $j<i$ olmak üzere önceki !!a{formula} $!B_j$'den zorunlulaştırma ile
  çıkan bir formül.
\end{enumerate}
$!B_n\ident !A$ olan böyle !!a{derivation} varsa, $!A$'nın
\emph{$\Sigma$'dan !!{derivable}} olduğunu söyler ve $\Sigma\Proves !A$
yazarız.
\end{defn}

Bu tanımla !!{derivable} formüller kümesinin normal bir modal mantık
oluşturduğu ve her !!{derivable} !!{formula} için, her aksiyomun doğru olduğu
her modelde doğruluk elde edildiği ortaya çıkacaktır. !!{derivation}s için bu özelliğe
\emph{sağlamlık} denir. Tersi olan \emph{tamlığı} ispatlamak daha zordur.

\end{document}
